\chapter{Bias Management --- RL's Real Game}
\label{ch:bias-management}

Newcomers to RL often assume the goal is to find an unbiased estimator of the
optimal policy.  This is a misunderstanding.  Every component of the RL
pipeline---sampling, estimation, target construction---introduces bias.  The
real engineering question is not ``how do I eliminate bias?'' but rather
\emph{which biases am I willing to accept, and how do I keep them under
control?}

\begin{keyidea}[The Bias Management Principle]
RL does not pursue unbiased optimization.  It selects, among a landscape of
unavoidable biases, the bias structure that is most stable and controllable in
practice.
\end{keyidea}


%% ================================================================
\section{Seven Sources of Bias}
\label{sec:seven-biases}

Almost every design choice in RL introduces a distinct form of bias.
Table~\ref{tab:bias-sources} catalogs the seven most important ones.

\begin{table}[ht]
\centering
\caption{Seven sources of bias in reinforcement learning.}
\label{tab:bias-sources}
\begin{tabular}{@{}lp{8.5cm}@{}}
\toprule
\textbf{Bias source} & \textbf{Mechanism} \\
\midrule
Sampling bias &
  The agent only visits states and actions that its current (or past) policy
  selects.  Regions of the state--action space that are never sampled can never
  be evaluated. \\
Bootstrapping bias &
  TD targets use the model's own estimates as labels:
  $y = r + \gamma \hat{Q}(s',a')$.  Errors in $\hat{Q}$ feed back into the
  next round of fitting. \\
Max / argmax bias &
  $\mathbb{E}[\max_a \hat{Q}(s,a)] \geq \max_a \mathbb{E}[\hat{Q}(s,a)]$.
  Taking the max over noisy estimates systematically overestimates value. \\
Function approximation bias &
  Neural networks generalize---and therefore \emph{mis}-generalize.  A network
  that fits well in visited regions may produce arbitrary values elsewhere. \\
Replay distribution bias &
  Off-policy methods reuse transitions sampled by old policies.  The
  distribution in the buffer drifts away from the current policy's visitation
  distribution. \\
Reward design bias &
  Shaped rewards, sparse rewards, and proxy rewards all inject the designer's
  assumptions.  The agent optimizes the reward it is given, not the reward the
  designer intended. \\
Policy-induced data bias &
  The policy generates its own training data.  A prematurely peaked policy
  narrows the data distribution, making it harder for the estimator to
  self-correct. \\
\bottomrule
\end{tabular}
\end{table}

These biases are not independent---they interact.  Bootstrapping amplifies
function-approximation error; max bias turns that amplified error into
overestimation; overestimation warps the policy, which in turn narrows sampling.
Understanding these feedback loops is the key to understanding why RL algorithms
look the way they do.


%% ================================================================
\section{Q-Learning's Bias: Max Overestimation}
\label{sec:q-bias}

Q-learning constructs its target as
\begin{equation}
  y = r + \gamma \max_{a'} Q_{\text{target}}(s', a').
\end{equation}
If each $\hat{Q}(s',a')$ has zero-mean noise $\epsilon_{a'}$, then
\begin{equation}
  \mathbb{E}\!\bigl[\max_{a'} \hat{Q}(s',a')\bigr]
  \;>\; \max_{a'}\mathbb{E}\!\bigl[\hat{Q}(s',a')\bigr],
\end{equation}
because the max operator preferentially selects actions whose noise happens to
be positive.  This is Jensen's inequality applied to the convex $\max$
function.

Over many bootstrapping steps, the overestimation compounds.  The agent begins
to ``believe'' that certain state--action pairs are far more valuable than they
actually are, and the policy chases these phantoms.

\subsection{Correction 1: Double Q-Learning}

Double DQN decouples action selection from value evaluation:
\begin{equation}
  a^{\ast} = \operatorname{argmax}_{a'} Q_{\text{online}}(s', a'),
  \qquad
  y = r + \gamma\, Q_{\text{target}}(s', a^{\ast}).
\end{equation}
The online network picks the action; the target network evaluates it.  Because
two different networks are unlikely to share the same noise pattern,
overestimation is reduced (though not eliminated).

\subsection{Correction 2: Target Network}

A target network $Q_{\text{target}}$ is a slow-moving copy of $Q_{\text{online}}$.
By not chasing every gradient step, it stabilizes the bootstrapping target and
dampens the positive-feedback loop between overestimation and policy shift.

\subsection{Correction 3: Clipped Double Q / Conservative Q}

TD3 maintains two critics and takes the minimum:
\begin{equation}
  y = r + \gamma\, \min\!\bigl(Q_1(s', \mu_{\text{target}}(s')),\;
  Q_2(s', \mu_{\text{target}}(s'))\bigr).
\end{equation}
This deliberately introduces a \emph{pessimistic} bias to counteract the
optimistic bias of $\max$.  The trade-off is slower value propagation, but
the critic surface the actor exploits is far more trustworthy.

Conservative Q-Learning (CQL) pushes the idea further in the offline setting by
adding a penalty that explicitly lowers $Q$ values for out-of-distribution
actions, ensuring the policy does not exploit regions where the critic has no
data support.

\begin{keyidea}[The Q-Learning Bias Correction Arc]
$\max$ overestimation $\;\longrightarrow\;$
Double Q (decouple selection/evaluation) $\;\longrightarrow\;$
Target net (slow the bootstrap) $\;\longrightarrow\;$
Clipped double Q (pessimistic bias) $\;\longrightarrow\;$
Conservative Q (penalize unsupported actions).
Each step trades a different form of controlled bias for a reduction in
uncontrolled overestimation.
\end{keyidea}


%% ================================================================
\section{Policy Gradient's Bias: Variance That Masquerades as Bias}
\label{sec:pg-bias}

Policy gradient methods avoid the $\max$ operator entirely.  The gradient
\begin{equation}
  \nabla J(\theta) = \mathbb{E}_{s,a \sim \pi_\theta}
  \bigl[\nabla \log \pi_\theta(a|s)\; A^{\pi}(s,a)\bigr]
\end{equation}
is, in principle, \emph{unbiased}---provided the advantage $A^{\pi}$ is computed
exactly and the actions are sampled from the current policy.

In practice, neither condition holds.  The advantage is estimated from a finite
number of noisy rollouts, and Monte Carlo returns have high variance.  High
variance means that the gradient signal is dominated by noise, and training
becomes unstable or painfully slow.

\subsection{Baseline: Centering the Advantage}

Subtracting a baseline $b(s)$ from the return does not change the expected
gradient (the $\nabla\log\pi$ term integrates to zero under the policy
distribution), but it dramatically reduces variance:
\begin{equation}
  \nabla J(\theta) = \mathbb{E}\bigl[\nabla \log \pi_\theta(a|s)\;
  \bigl(A(s,a) - b(s)\bigr)\bigr].
\end{equation}
The value function $V(s)$ is the most common baseline, yielding the advantage
$A(s,a) = Q(s,a) - V(s)$.

\subsection{GAE: Bias--Variance Trade-off in Advantage Estimation}

Generalized Advantage Estimation (GAE) interpolates between low-bias,
high-variance Monte Carlo returns and high-bias, low-variance one-step TD
errors via the parameter $\lambda$:
\begin{equation}
  \hat{A}_t^{\text{GAE}(\gamma,\lambda)}
  = \sum_{l=0}^{\infty} (\gamma\lambda)^l\, \delta_{t+l},
  \qquad
  \delta_t = r_t + \gamma V(s_{t+1}) - V(s_t).
\end{equation}
When $\lambda = 1$, GAE reduces to the full Monte Carlo advantage (unbiased,
high variance).  When $\lambda = 0$, it reduces to the one-step TD advantage
(biased, low variance).  Practitioners typically choose $\lambda \in [0.9, 0.99]$,
accepting a small amount of bias in exchange for much lower variance.

\subsection{PPO Clipping: Trading Bias for Stability}

PPO uses the importance-sampled objective
\begin{equation}
  L^{\text{CLIP}}(\theta) = \hat{\mathbb{E}}\!\left[\min\!\left(
    r_t(\theta)\,\hat{A}_t,\;
    \operatorname{clip}\!\bigl(r_t(\theta),\, 1{-}\epsilon,\, 1{+}\epsilon\bigr)\,
    \hat{A}_t
  \right)\right],
\end{equation}
where $r_t(\theta) = \pi_\theta(a_t|s_t)\,/\,\pi_{\text{old}}(a_t|s_t)$.

The clipping deliberately prevents the ratio from moving too far from 1.  This
means the gradient is no longer the true policy gradient---it is biased.  But
the bias is bounded and predictable, whereas the unclipped gradient can produce
catastrophic policy updates when the importance ratio explodes.

\begin{keyidea}[PPO's Bargain]
PPO intentionally introduces gradient bias through clipping.  The payoff is
engineering stability: updates that are ``close enough'' to the true gradient but
never catastrophically wrong.  This is perhaps the clearest example in RL of
\emph{choosing a more controllable bias structure}.
\end{keyidea}


%% ================================================================
\section{SAC's Bias: Entropy as Intentional Objective Rewriting}
\label{sec:sac-bias}

Soft Actor-Critic optimizes a modified objective:
\begin{equation}
  J_{\text{SAC}}(\pi) = \sum_t \mathbb{E}\!\bigl[
    r(s_t, a_t) + \alpha\,\mathcal{H}\!\bigl(\pi(\cdot|s_t)\bigr)
  \bigr],
\end{equation}
where $\alpha$ is a temperature parameter and $\mathcal{H}$ is the entropy of
the policy.

This is \emph{not} the original reward-maximization objective.  SAC admits an
intentional bias: the optimal SAC policy is not the same as the optimal
reward-maximizing policy.  It is the optimal \emph{entropy-regularized} policy.

Why accept this bias?  Three reasons:
\begin{enumerate}[nosep]
  \item \textbf{Stable exploration.}  The entropy term prevents the policy from
    collapsing onto a single action.  This keeps the sampling distribution broad
    and the critic's training data diverse.
  \item \textbf{Smoother Q landscape.}  When the policy stays stochastic, the
    Bellman target integrates over actions rather than hinging on a single
    $\max$.  This suppresses the spike-chasing behavior that plagues
    deterministic actor-critics.
  \item \textbf{Robustness.}  A stochastic policy is inherently more robust to
    perturbations---it hedges rather than committing to a single action.
\end{enumerate}

SAC's bias structure can be summarized as:

\begin{keyidea}[SAC's Entropy Bias]
SAC trades a soft objective (entropy-regularized reward) for stable exploration
and better data coverage.  The bias is explicit, tunable via $\alpha$, and
produces policies that are both effective and robust.
\end{keyidea}


%% ================================================================
\section{Choosing a Controllable Bias Structure}
\label{sec:choosing-bias}

Table~\ref{tab:bias-strategies} summarizes how major algorithm families handle
the bias problem.

\begin{table}[ht]
\centering
\caption{Bias strategies across algorithm families.}
\label{tab:bias-strategies}
\begin{tabular}{@{}lp{4cm}p{4.5cm}@{}}
\toprule
\textbf{Algorithm family} & \textbf{Dominant bias} & \textbf{Mitigation strategy} \\
\midrule
Q-learning / DQN &
  Max overestimation, bootstrap &
  Double Q, target net, conservative Q \\
DDPG / TD3 &
  Actor exploits critic errors &
  Clipped double critic, delayed update, target smoothing \\
REINFORCE / A2C &
  High variance (effective bias in finite samples) &
  Baseline, GAE, advantage normalization \\
PPO &
  Clipping introduces gradient bias &
  Bounded, predictable bias traded for stability \\
SAC &
  Entropy rewrites the objective &
  Intentional bias for exploration; tunable via $\alpha$ \\
\bottomrule
\end{tabular}
\end{table}

The pattern is clear.  No algorithm eliminates bias.  Each algorithm
\emph{chooses} which biases to accept based on what is controllable in its
particular setting:

\begin{itemize}[nosep]
  \item Value-based methods accept bootstrapping and approximation bias, then
    spend engineering effort suppressing max overestimation.
  \item Policy gradient methods are unbiased in theory but high-variance in
    practice; they introduce controlled bias (baselines, clipping, GAE's
    $\lambda$) to tame variance.
  \item SAC openly rewrites the objective, accepting an entropy bias that is
    explicit and tunable.
\end{itemize}

\begin{keyidea}[The Meta-Lesson]
The history of RL algorithm design is not a march toward unbiased estimation.
It is a search for \textbf{better bias structures}---ones whose distortions are
predictable, bounded, and tunable, even if the resulting objective is no longer
the ``pure'' reward-maximization problem.
\end{keyidea}


%% ================================================================
\section{Why This Matters in Practice}
\label{sec:bias-practice}

When debugging an RL system, the first question should not be ``is my
implementation correct?'' but rather ``which bias is dominating, and is it the
one I designed for?''

A few practical heuristics:

\begin{enumerate}[nosep]
  \item \textbf{If Q values diverge upward}: max overestimation is likely
    compounding.  Switch to Double Q or clipped double critics.
  \item \textbf{If the policy collapses to a single action}: sampling bias has
    taken over.  Increase entropy regularization or exploration noise.
  \item \textbf{If training is stable but performance plateaus}: the bias
    structure may be too conservative.  Try reducing PPO's clip range, lowering
    SAC's $\alpha$, or using less aggressive target-network lag.
  \item \textbf{If off-policy training is unstable}: replay distribution bias
    may be severe.  Consider prioritized replay, shorter buffer horizons, or
    switching to a near-on-policy method.
\end{enumerate}

Understanding bias is not an academic exercise---it is the most practical skill
in RL engineering.  Every hyperparameter in a modern RL algorithm is, at its
core, a knob that adjusts the bias--variance--stability trade-off.
